% Section 4: The Constitutional Enclosure Theorem
% IEEE Computer — "Silent Decisions Are Type Errors"
% Authors: Soares et al., 2026

\section{The Constitutional Enclosure Theorem}
\label{sec:theorem}

This section formalises the central claim of the paper: in any Safe Rust
program that uses the \texttt{silent-decisions-proof} crate, producing a
\emph{silent decision} is a compile-time impossibility.  We state the
theorem through three definitions, three axioms, and a constructive proof
with exhaustive case analysis.

% ----------------------------------------------------------------
\subsection{Definitions}
\label{sec:definitions}

\begin{definition}[Materialized Verdict]
\label{def:materialized}
A verdict $V$ is \emph{materialized} if and only if there exist an
evidence token $E$, a compliance token $C$, a decision outcome
$d \in \{\mathit{Allow},\, \mathit{Deny}\}$, and an explanation string
$x$ such that
\[
  V = \texttt{Verdict::new}(E,\, C,\, d,\, x).
\]
The function \texttt{Verdict::new} is the \emph{sole} constructor of
the type \texttt{Verdict}; no other term of that type is well-typed in
Safe Rust.
\end{definition}

\begin{definition}[Silent Decision]
\label{def:silent}
Let $\mathcal{S}$ denote the set of \emph{silent decisions}---decision
acts that carry legal consequences but lack a non-repudiable evidence
chain.  Formally, a decision act $\delta$ belongs to $\mathcal{S}$ if
and only if $\delta$ cannot be expressed as a Materialized Verdict
(Definition~\ref{def:materialized}).
\end{definition}

% ----------------------------------------------------------------
\subsection{Axioms}
\label{sec:axioms}

\begin{axiom}[Resource Introduction, $\otimes$-I]
\label{ax:intro}
If $\Gamma \vdash E : \texttt{EvidenceToken}$ and
$\Delta \vdash C : \texttt{ComplianceToken}$, then
\[
  \Gamma,\, \Delta \;\vdash\; \texttt{Verdict::new}(E,\, C,\, d,\, x)
  : \texttt{Verdict}.
\]
The contexts $\Gamma$ and $\Delta$ are \emph{disjoint}: linear logic
forbids sharing a resource between both premises.  In Rust this is
enforced by the move semantics of \texttt{Verdict::new}, which takes
$E$ and $C$ by value.  Once passed to the constructor, neither token
is accessible to the caller.
\end{axiom}

\begin{axiom}[Resource Uniqueness and Anti-Weakening]
\label{ax:unique}
\texttt{EvidenceToken} is annotated \texttt{\#[must\_use]} and does
not implement \texttt{Clone} or \texttt{Copy}.  Consequently:
\begin{enumerate}
  \item \emph{Contraction} is prohibited: a token cannot be
    duplicated, preventing the same evidence from being
    associated with more than one verdict.
  \item \emph{Weakening} is prohibited: a token cannot be silently
    discarded.  The compiler emits a \texttt{unused\_must\_use}
    diagnostic---treated as an error in this crate---whenever an
    \texttt{EvidenceToken} value is dropped without being consumed.
\end{enumerate}
The method \texttt{consume()} is the unique destructor; it transfers
the inner \texttt{Blake3Hash} to the \texttt{Verdict} constructor and
is declared \texttt{pub(crate)}, making it inaccessible outside the
defining crate.
\end{axiom}

\begin{axiom}[Encapsulation Boundary]
\label{ax:encap}
Let $\Gamma_{\mathit{ext}}$ denote any typing context external to the
defining crate (i.e., any context in a crate that depends on
\texttt{silent-decisions-proof} but is not the crate itself).
The following are ill-typed in $\Gamma_{\mathit{ext}}$:
\begin{enumerate}
  \item Struct-literal construction of \texttt{Verdict}---all fields
    are private (compiler error \texttt{E0451}).
  \item Direct invocation of \texttt{EvidenceToken::consume()}---the
    method is \texttt{pub(crate)} (compiler error \texttt{E0624}).
  \item Construction of \texttt{Blake3Hash} via its tuple-struct
    field---the field is private (compiler error \texttt{E0423}).
\end{enumerate}
Together, these restrictions close the encapsulation perimeter: the
only well-typed path to a \texttt{Verdict} in $\Gamma_{\mathit{ext}}$
is through \texttt{Verdict::new}, which enforces
Axiom~\ref{ax:intro}.
\end{axiom}

% ----------------------------------------------------------------
\subsection{Threat Model}
\label{sec:threat}

The theorem holds under the \emph{Safe Rust} memory model.
The \texttt{unsafe} keyword grants the programmer the ability to
bypass borrow-checker and visibility rules; consequently, code that
uses \texttt{unsafe} is explicitly out of scope.

This is not a practical limitation in governance contexts: production
deployments can enforce the absence of \texttt{unsafe} code via
\texttt{cargo geiger}~\cite{cargogeiger}, which reports every
\texttt{unsafe} block transitively reachable from the dependency
graph.  A zero-\texttt{unsafe} audit guarantees that the theorem
applies to the entire compiled binary.

Macro-generated code is in scope: procedural macros run before the
type checker and cannot produce terms that circumvent visibility rules
without themselves using \texttt{unsafe}.

% ----------------------------------------------------------------
\subsection{Proof of the Constitutional Enclosure Theorem}
\label{sec:proof}

\begin{theorem}[Constitutional Enclosure]
\label{thm:enclosure}
In any Safe Rust program $P$ whose crate graph includes
\texttt{silent-decisions-proof}, no term of type \texttt{Verdict}
can be constructed without consuming exactly one \texttt{EvidenceToken}
and exactly one \texttt{ComplianceToken}.  Equivalently,
$\delta \in \mathcal{S} \Rightarrow P$ does not type-check.
\end{theorem}

\begin{proof}
We show that every possible syntactic strategy for constructing a
\texttt{Verdict} without going through \texttt{Verdict::new} is
ill-typed.  The cases are exhaustive.

\medskip
\noindent\textbf{Case~A: Struct-literal construction.}
An adversary attempts to write
\begin{verbatim}
  Verdict { decision: Decision::Allow, .. }
\end{verbatim}
from a crate external to \texttt{silent-decisions-proof}.
By Axiom~\ref{ax:encap}(1), all fields of \texttt{Verdict} are
private.  The compiler rejects this expression with
\texttt{error[E0451]: field `evidence\_id` of struct `Verdict` is
private}.  This is verified by the compile-fail test
\texttt{tests/ui/verdict\_struct\_literal.rs}.

\medskip
\noindent\textbf{Case~B: External invocation of \texttt{consume()}.}
An adversary attempts to call
\texttt{EvidenceToken::consume()} directly from an external crate,
bypassing \texttt{Verdict::new} to extract a \texttt{Blake3Hash} and
assemble a \texttt{Verdict} fragment.
By Axiom~\ref{ax:encap}(2), \texttt{consume()} is declared
\texttt{pub(crate)}.  The compiler rejects the call with
\texttt{error[E0624]: method `consume` is private}.
This is verified by the compile-fail test
\texttt{tests/ui/external\_consume\_call.rs}.

\medskip
\noindent\textbf{Case~C: Silent drop of \texttt{EvidenceToken}.}
A developer issues a decision without providing evidence by
constructing an \texttt{EvidenceToken} and immediately dropping it
(or ignoring the return value of a function that returns one).
By Axiom~\ref{ax:unique}, \texttt{EvidenceToken} is \texttt{\#[must\_use]}.
The compiler emits \texttt{warning: unused return value of
`EvidenceToken` that must be used}, and because the crate sets
\texttt{\#![deny(unused\_must\_use)]}, this warning is promoted to a
hard error.  Note that \texttt{let \_ = expr} in Rust explicitly silences the
\texttt{\#[must\_use]} lint without consuming the token.  This does
not violate the theorem: no \texttt{Verdict} is produced, so no
silent decision (Definition~\ref{def:silent}) materializes.  The
\texttt{\#[must\_use]} annotation serves as a developer-facing
guardrail against accidental drops, not as a formal closure of the
type perimeter---that role belongs to Axiom~\ref{ax:encap} (private
fields and \texttt{pub(crate)} visibility).  The developer cannot
construct a \texttt{Verdict} without going through
\texttt{Verdict::new}.
This is verified by the compile-fail test
\texttt{tests/ui/dropped\_evidence\_token.rs}.

\medskip
Since Cases~A, B, and C are exhaustive over all syntactic strategies
for circumventing the constructor, and all three are ill-typed in
Safe Rust, no term of type \texttt{Verdict} can be produced in
$\Gamma_{\mathit{ext}}$ without passing through
\texttt{Verdict::new}.  By Definition~\ref{def:materialized}, every
such term is a Materialized Verdict, and by
Definition~\ref{def:silent}, no silent decision $\delta \in
\mathcal{S}$ is representable.
\end{proof}

% ----------------------------------------------------------------
\subsection{Corollary: Fail-Secure by Construction}
\label{sec:corollary}

\begin{corollary}[Fail-Secure]
\label{cor:failsecure}
In any Safe Rust program satisfying the conditions of
Theorem~\ref{thm:enclosure}, the code path that handles a runtime
failure to produce evidence (e.g., a database timeout, a missing
context field) cannot silently issue a \texttt{Verdict}---the type
system ensures at compile time that every branch leading to a
\texttt{Verdict} must supply a valid \texttt{EvidenceToken}.
The system \emph{fails securely}: it cannot issue a decision that
lacks an evidence chain.
\end{corollary}

\begin{proof}
If the evidence-production path is absent or returns an error type
(e.g., \texttt{Result<EvidenceToken, E>}), the only well-typed
continuation that produces a \texttt{Verdict} requires unwrapping
the result and handling the error branch explicitly.  The error
branch cannot produce a \texttt{Verdict} without evidence; it must
either propagate the error upward or halt.  By
Theorem~\ref{thm:enclosure}, any attempt to construct a
\texttt{Verdict} in the error branch without a valid
\texttt{EvidenceToken} is ill-typed.
\end{proof}

The Fail-Secure Corollary aligns with the regulatory intent of
LGPD Art.~18\S2 and EU AI Act Art.~86: the obligation to provide an
explanation is not advisory but \emph{structural}---it cannot be
omitted under load, behind a feature flag, or through a code path
that bypasses an \texttt{assert}.
